% 来源：OEBPS/Text/01_53.xhtml（OPF spine 顺序；仅第一册）
\cleardoublepage
\BookEntry[安迪·希尔《要快，但别赶！》]{安迪·希尔《要快，但别赶！》}

\BookCover{images/书单36要快，但别赶！.jpg}{[美] 安迪·希尔、约翰·伍登著| Simon \& Schuster 2001年版}

\begin{BookSummaryBox}
\textbf{世界上最优秀的篮球教练之一伍登的训练法则}
\par
\end{BookSummaryBox}

\subsubsection*{这本书谈什么？}

1968年，安迪·希尔是洛杉矶市大学高级中学的出色后卫。他平均每场比赛能够得到27分，许多篮球强校向他伸出了橄榄枝。他选择了加州大学洛杉矶分校，在大一的队伍中打出了极佳的表现。他平均每场比赛能拿到19分和8个助攻，和亨利·毕比一起被评为球队最有价值球员。
\par

然而，安迪在加州大学洛杉矶分校的灰姑娘故事很快就结束了。该校校队刚刚获得七年内的第六个全国冠军，希尔很快就发现他自己达不到伍登教练对首发球员的要求。和其他在大学校队期间坐板凳的替补队员一样，在伍登教练手下，希尔也深刻地体会到体育运动的非平等主义本质。在他这本令人愉快的《要快，但别赶！》中，希尔叙述了那段痛苦经历、他和伍登最后的和解，以及他从这段经历中学到的宝贵经验。
\par

\subsubsection*{芒格为什么谈起这本书？}

芒格在南加州大学GOULD法学院毕业典礼上的演讲中，谈到一个他认为很重要的道理，那就是：将不平等最大化通常能收到奇效。这句话是什么意思呢？芒格是这样说的：
\par

\begin{quotation}
伍登曾经是世界上最优秀的篮球教练。他对五个水平较低的球员说：“你们不会得到上场的时间——你们是陪练。”比赛几乎都是那七个水平较高的球员在打的。嗯，这七个水平高的球员学到了更多——别忘了学习机器的重要性——因为他们独享了所有的比赛时间。在他采用非平等主义的方法时，伍登比从前赢得了更多的比赛。我认为生活就像比赛，也充满了竞争，我们要让那些最有能力和最愿意成为学习机器的人发挥最大的作用。
\par
\end{quotation}

\subsubsection*{《要快，但别赶！》作者安迪·希尔如是说}

“拥有最优秀球员的球队几乎总是会赢。”
\par

“要快，但别赶。万事确定的时刻永远不会到来，但是你必须一直保持平衡。”
\par

“如果在准备这件事上失败，那就是在为失败做准备。”
\par

“取胜的球队是那些失误最多的球队。”
\par
